\documentclass[11pt, a4paper]{article}
\usepackage{amsthm}
\usepackage{amsmath}
\usepackage{amssymb}
\usepackage{geometry}
\geometry{
    margin=1in,
    headheight=20pt,
    footskip=10mm
}
\title{3}
\author{Nagusa, G}
\date{\today}

\begin{document}
\maketitle

\section*{Problem}
Let $R$ be a commutative ring.
We say that an element $a \in R$ is nilpotent if
\[
a^n = 0_R \text{ for some positive integer } n
\]
Prove that the subset of $R$ consisting of all nilpotent elements is a subring of $R$.

\section*{Proof}
Let $R$ be a commutative ring, and let $S$ be the set of all nilpotent elements in $R$:
\[ S = \{ r \in R \mid r^n = 0_R \text{ for some positive integer } n \} \] 
We will verify the following to prove that $S$ is a subring of $R$:
\begin{enumerate}
    \item $0_R \in S$
    \item $a + (-b) \in S$ for all $a, b \in S$
    \item $a \cdot b \in S$ for all $a, b \in S$
\end{enumerate}

\subsection*{1.}
We observe that ${0_R}^1 = 0_R$. \\
Thus, $0_R$ is nilpotent, which means $0_R \in S$.

\subsection*{2.}
We will prove this in two steps$((-a) \in S$ and $a+b \in S)$.
\begin{enumerate}
    \item $(-a) \in S$\\[5pt]
    Let $a \in S$. Then there is some positive integer $n$ such that $a^n = 0_R$. \\
    Since $R$ is a ring, $(-a) \in R$.
    \begin{itemize}
        \item[--] \textbf{Case 1: $n$ is even.} \\
        Let $n = 2k$ for some integer $k$. \\
        Then $(-a)^n = (-a)^{2k} = ((-a)^2)^k = (a^2)^k = a^{2k} = a^n = 0_R$.
        \item[--] \textbf{Case 2: $n$ is odd.} \\
        Let $n = 2k+1$ for some integer $k$. \\
        Then $(-a)^n = (-a)^{2k+1} = (-a)^{2k}\cdot(-a) = a^{2k}\cdot(-a) = -(a^{2k}\cdot a) = -a^{2k+1} = -a^n = -0_R = 0_R$.
    \end{itemize}
    Thus, $(-a)$ is nilpotent, which means $(-a) \in S$.
    
    \item $a+b \in S$\\[5pt]
    Let $a, b \in S$. Then there are some positive integers $n, m$ such that $a^n = 0_R$ and $b^m = 0_R$. \\
    Consider the expansion
    \[
    (a+b)^{n+m} = \underbrace{(a+b)(a+b)\cdots(a+b)}_{n+m \text{ times}}
    \]
    Since $R$ is a commutative ring, we can rearrange the multiplication order of $a$ and $b$ in each term. Thus, the expansion consists of a sum of terms of the form $a^i b^j$, where $i+j = n+m$.
    \[
    (a+b)^{n+m} = \sum (\text{terms of the form } a^i b^j)
    \]
    For each term $a^i b^j$, either $i \ge n$ or $j \ge m$ must hold (otherwise $i+j < n+m$).
    \begin{itemize}
        \item[--] \textbf{Case 1: $i \ge n$} \\
        $a^ib^j = (a^n \cdot a^{i-n})b^j = (0_R \cdot a^{i-n})b^j = 0_R$
        \item[--] \textbf{Case 2: $j \ge m$} \\
        $a^ib^j = a^i(b^m \cdot b^{j-m}) = a^i(0_R \cdot b^{j-m}) = 0_R$
    \end{itemize}
    Since every term $a^ib^j = 0_R$, we obtain
    \[
    (a+b)^{n+m} = 0_R
    \]
    Thus, $a+b$ is nilpotent, so $a+b \in S$.
\end{enumerate}
By steps 1 and 2, we have $a + (-b) \in S$ for all $a, b \in S$.

\subsection*{3.}
Let $a, b \in S$. Then there is some positive integer $n$ such that $a^n = 0_R$. \\
Since $R$ is a commutative ring,
\[
(a \cdot b)^n = a^n \cdot b^n = 0_R \cdot b^n = 0_R
\]
Thus, $a \cdot b$ is nilpotent, so $a \cdot b \in S$.

\subsection*{$\therefore$ $S$ is a subring of $R$}

\end{document}